% Luc Destombe --- licence CC BY-NC-SA 4.0
% https://creativecommons.org/licenses/by-nc-sa/4.0/deed.fr
% Fichier autonome : compiler avec pdflatex, deux fois (signets, nombre de pages).
\documentclass[a4paper,10pt]{article}
\makeatletter
\newif\iflycee@niveaux
\lycee@niveauxtrue
\newif\iflycee@palatino
\lycee@palatinotrue

\RequirePackage[utf8]{inputenc}
\RequirePackage[T1]{fontenc}
\RequirePackage[french]{babel}
\RequirePackage{etoolbox}

\newcommand{\ShorthandsOff}{\@ifundefined{shorthandoff}{}{\shorthandoff{;:!?}}}
\newcommand{\ShorthandsOn}{\@ifundefined{shorthandon}{}{\shorthandon{;:!?}}}
\ShorthandsOff
\AtBeginDocument{\ShorthandsOn}
\AtBeginEnvironment{tikzpicture}{\ShorthandsOff}

\RequirePackage{geometry}
\RequirePackage{fancyhdr}
\RequirePackage{lastpage}
\RequirePackage{multicol}
\RequirePackage{array}
\RequirePackage{enumitem}
\RequirePackage{xcolor}
\RequirePackage{needspace}

\RequirePackage{amsmath,amssymb}
\RequirePackage{mathpazo}
\DeclareMathSizes{10.95}{10.95}{8}{5.5}
\begingroup\catcode`\_=\active \catcode`\^=\active
\gdef_#1{\sb{\scriptscriptstyle #1}}
\gdef^#1{\sp{\scriptscriptstyle\lycee@moinsepais #1}}
\endgroup
\newcommand\lycee@moins{\mathbin{%
  \setbox\z@\hbox{$\m@th\scriptscriptstyle\lycee@moinsorig$}%
  \dimen@=.55\wd\z@ \advance\dimen@-.5pt
  \hbox to.75\wd\z@{\hss$\m@th\scriptscriptstyle\vcenter{\hbox{%
    \tikz\draw[line width=.5pt,line cap=round](0,0)--(\the\dimen@,0);}}$\hss}}}
\begingroup\lccode`\~=`\- \lowercase{\endgroup\let~}\lycee@moins
\newcommand\lycee@moinsepais{\mathcode`\-="8000 }
\AtBeginDocument{\mathchardef\lycee@moinsorig=\mathcode`\-\relax}
\AtBeginDocument{\catcode`\_=12 \mathcode`\_="8000
  \catcode`\^=12 \mathcode`\^="8000 }
\newcommand\lycee@exposants{%
  \fontdimen13\textfont2=.48\fontdimen6\textfont2
  \fontdimen14\textfont2=.48\fontdimen6\textfont2
  \fontdimen15\textfont2=.48\fontdimen6\textfont2
  \fontdimen13\scriptfont2=.48\fontdimen6\scriptfont2
  \fontdimen14\scriptfont2=.48\fontdimen6\scriptfont2
  \fontdimen15\scriptfont2=.48\fontdimen6\scriptfont2}
\every@math@size\expandafter{\the\every@math@size\lycee@exposants}
\newlength{\retraitcalculs}
\setlength{\retraitcalculs}{2em}

\RequirePackage{tikz}
\usetikzlibrary{arrows.meta,calc,babel,shapes.misc}
\RequirePackage[most]{tcolorbox}
\tcbuselibrary{skins,breakable}

\definecolor{coulDef}{HTML}{1F4E79}
\definecolor{coulProp}{HTML}{7030A0}
\definecolor{coulRem}{HTML}{C55A11}
\definecolor{coulEx}{HTML}{2E7D32}
\definecolor{coulExo}{HTML}{B03A2E}
\definecolor{coulPart}{HTML}{1F4E79}
\definecolor{coulMeth}{HTML}{1B6E4F}
\definecolor{coulCourbe}{HTML}{0B5ED7}
\definecolor{coulCourbeB}{HTML}{C00000}
\definecolor{coulCourbeC}{HTML}{2E7D32}
\definecolor{coulGrille}{HTML}{8C8C8C}
\definecolor{coulRep}{HTML}{1B6E4F}
\definecolor{coulBar}{HTML}{2E7D32}

\newcommand{\R}{\mathbb{R}}
\newcommand{\Rs}{\mathbb{R}^{*}}
\renewcommand{\le}{\leqslant}
\renewcommand{\ge}{\geqslant}
\renewcommand{\approx}{\simeq}
\newcommand{\vect}[1]{\overrightarrow{#1}}
\newcommand{\norme}[1]{\left\lVert #1 \right\rVert}
\newcommand{\intervalle}[2]{\left[#1\,;#2\right]}
\RequirePackage{cancel}
\renewcommand{\CancelColor}{\color{coulExo}}
\newcommand{\fois}[1]{\times{\color{coulExo}#1}}
\newcommand{\barre}[1]{\cancel{#1}}

\tikzset{grille/.style={coulGrille,line width=0.4pt}}
\newcommand{\quadrillage}[4]{\draw[grille,step=1] (#1,#2) grid (#3,#4);}
\tikzset{
  croix/.style={cross out,draw,line width=0.6pt,inner sep=0pt,minimum size=4pt},
  croixdroite/.style={croix,rotate=45,line width=0.8pt},
}
\newcommand{\pt}[3]{\path (#1) node[croix]{} node[#2,font=\small,inner sep=2.5pt]{$#3$};}

\newcommand{\lycee@repere}[4]{%
  \draw[grille,step=1] (#1,#2) grid (#3,#4);
  \draw[-{Stealth[length=2mm]},thick] (#1,0)--({#3+0.4},0) node[below left=-1pt]{$x$};
  \draw[-{Stealth[length=2mm]},thick] (0,#2)--(0,{#4+0.4}) node[below left=-1pt]{$y$};
  \node[below left,font=\scriptsize,inner sep=1.5pt] at (0,0) {$0$};}
\newcommand{\lycee@gradx}[1]{\foreach \k/\l in {#1}{%
  \draw (\k,0.07)--(\k,-0.07) node[below,font=\scriptsize,inner sep=1.5pt]{$\l$};}}
\newcommand{\lycee@grady}[1]{\foreach \k/\l in {#1}{%
  \draw (0.07,\k)--(-0.07,\k) node[left,font=\scriptsize,inner sep=1.5pt]{$\l$};}}
\newcommand{\lycee@intff}[2]{\left[#1\,;#2\right]}
\newcommand{\lycee@intfo}[2]{\left[#1\,;#2\right[}
\newcommand{\lycee@intof}[2]{\left]#1\,;#2\right]}
\newcommand{\lycee@intoo}[2]{\left]#1\,;#2\right[}
\newcommand{\lycee@ens}[1]{\left\{#1\right\}}
\AtBeginDocument{%
  \providecommand{\repere}{\lycee@repere}%
  \providecommand{\gradx}{\lycee@gradx}%
  \providecommand{\grady}{\lycee@grady}%
  \providecommand{\Cf}{\mathcal{C}\sb{\scriptscriptstyle f}}%
  \providecommand{\Cg}{\mathcal{C}\sb{\scriptscriptstyle g}}%
  \providecommand{\intff}{\lycee@intff}%
  \providecommand{\intfo}{\lycee@intfo}%
  \providecommand{\intof}{\lycee@intof}%
  \providecommand{\intoo}{\lycee@intoo}%
  \providecommand{\ens}{\lycee@ens}}

\newlist{questions}{enumerate}{3}
\setlist[questions,1]{label=\arabic*\textdegree),leftmargin=*,itemsep=2pt,topsep=3pt}
\setlist[questions,2]{label=\alph*),leftmargin=*,itemsep=1pt,topsep=2pt}
\setlist[questions,3]{label=\roman*),leftmargin=*,itemsep=1pt,topsep=1pt}
\setlist[itemize]{label=\textbullet,leftmargin=*,itemsep=2pt,topsep=3pt}

\newcommand{\besoinplace}[1]{\par\penalty\@M\begingroup
  \dimen@=#1\relax\dimen@ii=\pagegoal\advance\dimen@ii-\pagetotal
  \ifdim\dimen@>\dimen@ii\ifdim\dimen@ii>\z@\vfil\fi\break\fi\endgroup}

\newcommand{\lycee@pied}{\thepage/\pageref{LastPage}}
\newcommand{\entete}[2]{%
  \pagestyle{fancy}\fancyhf{}%
  \fancyhead[L]{\footnotesize\color{coulPart}\textbf{#1}}%
  \fancyhead[R]{\footnotesize\color{coulPart}\textbf{#2}}%
  \fancyfoot[C]{\footnotesize\lycee@pied}%
  \renewcommand{\headrulewidth}{0.4pt}%
  \renewcommand{\headrule}{\hbox to\headwidth{%
    \color{coulPart!40}\leaders\hrule height \headrulewidth\hfill}}}

\RequirePackage{ccicons}
\newcommand{\lycee@licence}{{\scriptsize\color{gray}%
  \copyright\ Luc Destombe --- \ccbyncsa\ CC BY-NC-SA 4.0}}
\newif\iflycee@licence \lycee@licencetrue
\newcommand{\sanslicence}{\lycee@licencefalse}
\AtBeginDocument{\iflycee@licence\AddToHookNext{shipout/foreground}{%
  \put(0,0){\rlap{\kern\dimexpr1in+\hoffset+\oddsidemargin+\textwidth\relax
    \llap{\raisebox{-\dimexpr1in+\voffset+\topmargin+\headheight+\headsep
      +\textheight+\footskip\relax}{\lycee@licence}}}}}\fi}

\newcommand{\formulesserrees}{%
  \setlength{\abovedisplayskip}{3pt}\setlength{\belowdisplayskip}{3pt}%
  \setlength{\abovedisplayshortskip}{2pt}\setlength{\belowdisplayshortskip}{3pt}}

\setlength{\parindent}{0pt}

\RequirePackage[implicit=false,hidelinks,pdfencoding=auto,psdextra,bookmarksopen,
  bookmarksopenlevel=1,pdfcreator={},pdfproducer={}]{hyperref}
\RequirePackage{bookmark}
\hypersetup{pdfauthor={Luc Destombe},pdfkeywords={CC BY-NC-SA 4.0}}
\newcounter{lycee@signet}
\newcommand{\signet}[3][]{\stepcounter{lycee@signet}%
  \edef\lycee@dest{\if\relax\detokenize{#1}\relax
    signet.\arabic{lycee@signet}\else#1\fi}%
  \hypertarget{\lycee@dest}{}\bookmark[level=#2,dest=\lycee@dest]{#3}}
\pdfstringdefDisableCommands{%
  \def\R{R}\def\vect#1{#1}\def\Cf{Cf}\def\Cg{Cg}\def\fois#1{×#1}%
  \def\barre#1{#1}\def\intervalle#1#2{[#1;#2]}\def\intff#1#2{[#1;#2]}%
  \def\textsuperscript#1{#1}\def\dfrac#1#2{#1/#2}\def\frac#1#2{#1/#2}%
  \def\vec#1{#1⃗}}%
\RequirePackage[official]{eurosym}

\geometry{top=0.8cm,bottom=1.3cm,left=1cm,right=1cm,footskip=0.6cm}

\pagestyle{fancy}
\fancyhf{}
\renewcommand{\headrulewidth}{0pt}
\fancyfoot[C]{\footnotesize Page \thepage{} / \pageref{LastPage}}

\newcommand{\titrefiche}[2]{%
  \begin{center}
    \Large\textbf{#1}\\[2pt]
    \large\textbf{#2}\\[2pt]
  \end{center}
  \vspace{0.10cm}}

\setlength{\columnseprule}{0.4pt}
\setlength{\columnsep}{15pt}
\renewcommand{\columnseprulecolor}{\color{black!45}}
\setlength{\multicolsep}{2pt}

\newcounter{partiefiche}
\renewcommand{\thepartiefiche}{\Roman{partiefiche}}
\newif\ifapresbandeau
\newcommand{\lycee@bandeau}[1]{%
  \par\penalty-600
  \vspace{0.08cm}%
  \noindent\colorbox{coulPart}{%
    \parbox{\dimexpr\linewidth-2\fboxsep\relax}{%
      \color{white}\bfseries\raggedright\strut#1\strut}}%
  \nopagebreak\par\nopagebreak
  \vspace{0.06cm}\nopagebreak
  \apresbandeautrue}
\newcommand{\partiefiche}[1]{%
  \refstepcounter{partiefiche}%
  \lycee@bandeau{\signet{1}{\thepartiefiche{} --- #1}\thepartiefiche{} --- #1}}
\newcommand{\partiefichesansnum}[1]{\lycee@bandeau{\signet{1}{#1}#1}}

\newcounter{exo}
\newlength{\espaceexo}\setlength{\espaceexo}{7pt}
\newlength{\espacetitre}\setlength{\espacetitre}{2pt}
\newcommand{\exo}[1][\relax]{%
  \par
  \ifapresbandeau\apresbandeaufalse\else\penalty-150\addvspace{\espaceexo}\fi
  \refstepcounter{exo}%
  \noindent\signet[exercice-\arabic{exo}]{2}{Exercice \arabic{exo}}%
  \textbf{\color{coulExo}\underline{Exercice \theexo}}%
  \ifx\relax#1\relax\else\textbf{ : #1}\fi
  \par\nopagebreak\vspace{\espacetitre}\nopagebreak
  \ignorespaces}

\setlist[enumerate]{nosep,leftmargin=*,itemsep=2pt,topsep=2pt}
\setlist[itemize]{nosep,leftmargin=*,topsep=2pt}
\newcommand{\choix}[3]{\par\hspace*{1em}%
  \makebox[0.3\linewidth][l]{a) #1}\makebox[0.3\linewidth][l]{b) #2}c) #3}
\newcommand{\deux}[2]{\par\hspace*{0.6em}%
  \makebox[0.48\linewidth][l]{#1}#2}
\newcommand{\trois}[3]{\par\hspace*{0.6em}%
  \makebox[0.32\linewidth][l]{#1}\makebox[0.32\linewidth][l]{#2}#3}

  \renewcommand{\exo}[1][\relax]{%
    \ifapresbandeau\apresbandeaufalse\else\par\penalty-150\fi
    \refstepcounter{exo}%
    \vspace{0.05cm}%
    \noindent\signet[exercice-\arabic{exo}]{2}{Exercice \arabic{exo}}%
    \textbf{\color{coulExo}\underline{Exercice \theexo}}%
    \ifx\relax#1\relax\else\textbf{ : #1}\fi%
    \par\nopagebreak
    \ignorespaces}
  \newcommand{\rappel}[1]{%
    \par\vspace{0.06cm}\nopagebreak
    \noindent\fcolorbox{black!35}{black!5}{%
      \parbox{\dimexpr\linewidth-2\fboxsep-2\fboxrule\relax}{%
        \small\textbf{Rappel.} #1}}%
    \par\vspace{0.06cm}\nopagebreak}
  \setlist[enumerate]{nosep,leftmargin=*}
  \setlist[itemize]{nosep,leftmargin=*}
  \setlength{\multicolsep}{3pt plus 1pt minus 1pt}
  \setlength{\premulticols}{10pt}
  \setlength{\postmulticols}{10pt}
  \newlist{expr}{enumerate}{1}
  \setlist[expr]{label={\Alph*\ $=$},nosep,leftmargin=*,itemsep=0pt}

\renewenvironment{center}{\par\addvspace{3pt}\centering}{\par\addvspace{3pt}}
\formulesserrees
\AtBeginDocument{\formulesserrees}
\clubpenalty=10000
\widowpenalty=10000
\displaywidowpenalty=10000
\brokenpenalty=10000
\setlength{\parskip}{0pt}

\RequirePackage{ragged2e}
\setlength{\RaggedRightRightskip}{0pt plus 1fil}
\AtBeginDocument{\RaggedRight\binoppenalty=10000 \relpenalty=10000 }
\g@addto@macro\@parboxrestore{\RaggedRight}
\makeatother

\newcommand{\coord}[2]{\begin{pmatrix}#1\\#2\end{pmatrix}}

\ShorthandsOff
\tikzset{
  vecteur/.style={-{Stealth[length=2.6mm]},line width=1.1pt,coulCourbe},
  vecteurB/.style={-{Stealth[length=2.6mm]},line width=1.1pt,coulCourbeB},
  vecteurC/.style={-{Stealth[length=2.6mm]},line width=1.1pt,coulCourbeC},
  aide/.style={dashed,black!45,line width=0.6pt},
  fig/.style={line width=0.8pt,black!75},
  axe/.style={aide,-{Stealth[length=2.2mm]}},
  deplacement/.style={-{Stealth[length=2.6mm]},line width=1.1pt,black!75},
}
\ShorthandsOn

\newcommand{\solidetire}[4]{%
\begin{tikzpicture}[scale=0.7]
  \draw[fig,fill=black!6] (-0.7,-0.45) rectangle (0.7,0.45);
  \draw[axe] (-1.6,0) -- (6.2,0);
  \draw[deplacement] (0,0) -- node[pos=0.6,below,font=\footnotesize]{déplacement\,: #4} (5,0);
  \draw (1.5,0) arc (0:#1:1.5);
  \node[font=\footnotesize,anchor=west] at ({1.55*cos(#1/2)},{1.55*sin(#1/2)}) {$#1^{\circ}$};
  \draw[vecteur] (0,0) -- node[midway,above left,font=\footnotesize]{#2\,: #3} ({3.4*cos(#1)},{3.4*sin(#1)});
  \pt{0,0}{below left}{G}
\end{tikzpicture}}

\newcommand{\troisforces}[6]{%
\begin{tikzpicture}[scale=0.7]
  \draw[fig,fill=black!6] (-0.7,-0.45) rectangle (0.7,0.45);
  \draw[axe] (-3.2,0) -- (6.2,0);
  \draw[deplacement] (0,0) -- node[pos=0.6,below,font=\footnotesize]{déplacement\,: #6} (5,0);
  \draw (1.5,0) arc (0:#1:1.5);
  \node[font=\footnotesize,anchor=west] at ({1.55*cos(#1/2)},{1.55*sin(#1/2)}) {#2};
  \draw[vecteur] (0,0) -- node[midway,above left,font=\footnotesize]{#3} ({3.2*cos(#1)},{3.2*sin(#1)});
  \draw[vecteurB] (0,0) -- node[pos=0.8,right,font=\footnotesize]{#4} (0,-2.4);
  \draw[vecteurC] (0,0) -- node[pos=0.8,above,font=\footnotesize]{#5} (-2.6,0);
  \pt{0,0}{below left}{G}
\end{tikzpicture}}

\newcommand{\vecteuraxe}[1]{%
\begin{tikzpicture}[scale=0.55]
  \draw[axe] (-2.6,0) -- (4.2,0) node[right,font=\footnotesize]{axe};
  \draw[deplacement] (0,0) -- node[pos=0.8,below,font=\footnotesize]{$\vec{u}$} (2.5,0);
  \draw (1.1,0) arc (0:#1:1.1);
  \node[font=\footnotesize,anchor=west] at ({1.2*cos(#1/2)},{1.2*sin(#1/2)}) {$#1^{\circ}$};
  \draw[vecteur] (0,0) -- node[midway,left,font=\footnotesize]{$\vec{v}$} ({3*cos(#1)},{3*sin(#1)});
\end{tikzpicture}}

\begin{document}

\titrefiche{1\textsuperscript{re} STI2D --- Produit scalaire}{Fiche d'exercices du chapitre}

\begin{multicols*}{2}\raggedcolumns

\partiefiche{Coordonnées d'un vecteur}

\exo[Lire et tracer]
Dans le repère ci-dessous :
\begin{enumerate}[label=\alph*)]
  \item tracer les vecteurs $\vect{AB}$, $\vect{CD}$ et $\vect{EF}$, puis lire leurs
        coordonnées ;
  \item tracer les représentants d'origine $G$ des vecteurs $\vec{u}\coord{2}{3}$ et
        $\vec{v}\coord{-1}{-2}$.
\end{enumerate}
\begin{center}
\begin{tikzpicture}[x=0.45cm,y=0.45cm]
  \repere{-4}{-4}{5}{4}
  \gradx{-3,-2,-1,1,2,3,4} \grady{-3,-2,-1,1,2,3}
  \pt{-3,1}{above left}{A} \pt{1,3}{above}{B} \pt{2,-1}{below}{C} \pt{4,2}{above}{D}
  \pt{-2,-2}{below left}{E} \pt{1,-3}{below right}{F} \pt{-3,-1}{below left}{G}
\end{tikzpicture}
\end{center}

\exo[Calculer des coordonnées]
On donne les points $K(3\,;-1)$, $L(-2\,;4)$, $M(0\,;-5)$ et $N(-4\,;-3)$. Calculer les
coordonnées des vecteurs $\vect{KL}$, $\vect{LK}$, $\vect{MN}$, $\vect{NK}$ et $\vect{LM}$.

\exo[Relation de Chasles]
\rappel{relation de Chasles : $\vect{AB}+\vect{BC}=\vect{AC}$ (l'arrivée du premier
vecteur est le départ du second) ; $\vect{BA}=-\vect{AB}$.}
Écrire chaque expression avec un seul vecteur.
\begin{enumerate}[label=\alph*)]
  \item $\vect{AB}+\vect{BC}$
  \item $\vect{MN}+\vect{NE}+\vect{EF}$
  \item $\vect{CA}+\vect{AB}$
  \item $\vect{AB}+\vect{BA}$
  \item $\vect{AC}-\vect{BC}$
  \item $\vect{KL}+\vect{MK}$
\end{enumerate}

\exo[Le trajet d'un drone]
Un drone vole à hauteur constante. Le repère est en mètres : l'axe des abscisses est
dirigé vers l'est, celui des ordonnées vers le nord.
\begin{enumerate}[label=\alph*)]
  \item Le drone part de $D(2\,;1)$ et avance de $5$~m vers l'est et de $3$~m vers le
        sud. Donner les coordonnées du vecteur $\vec{u}$ de ce déplacement, puis du point
        d'arrivée $E$.
  \item Il rejoint ensuite le point $F(-1\,;4)$. Calculer les coordonnées de
        $\vect{EF}$, puis décrire ce déplacement par une phrase complète (« le drone
        avance de … vers … »).
\end{enumerate}

\partiefiche{Somme de vecteurs, produit par un nombre réel}

\exo[Calculer avec les coordonnées]
On donne $\vec{u}\coord{2}{-3}$, $\vec{v}\coord{-4}{1}$ et $\vec{w}\coord{0}{2}$. Calculer
les coordonnées des vecteurs :
\begin{enumerate}[label=\alph*)]
  \item $\vec{u}+\vec{v}$ ;
  \item $\vec{u}-\vec{v}$ ;
  \item $4\vec{w}$ ;
  \item $2\vec{u}+3\vec{v}$ ;
  \item $\vec{u}-2\vec{v}+\vec{w}$.
\end{enumerate}

\exo[Avec des points]
On donne $A(-1\,;2)$, $B(3\,;0)$ et $C(1\,;-4)$. Calculer les coordonnées des vecteurs
$\vect{AB}+\vect{AC}$, $2\vect{AB}$, $\vect{AB}-\vect{BC}$ et $3\vect{AC}-2\vect{AB}$.

\exo[Trouver un vecteur]
Pour chaque cas, calculer les coordonnées du vecteur inconnu, en posant ses
coordonnées $x$ et $y$.
\begin{enumerate}[label=\alph*)]
  \item $\vec{u}\coord{1}{-2}$ ; trouver $\vec{v}$ tel que $\vec{u}+\vec{v}\coord{3}{4}$.
  \item $\vec{a}\coord{-3}{5}$ ; trouver $\vec{b}$ tel que $\vec{a}+\vec{b}\coord{0}{-1}$.
  \item $\vec{u}\coord{1}{-2}$ ; trouver $\vec{w}$ tel que $2\vec{u}+\vec{w}\coord{5}{1}$.
\end{enumerate}

\exo[Deux remorqueurs]
Deux remorqueurs tirent un navire, placé en $O$, avec les forces $\vec{F_{1}}\coord{3}{2}$
et $\vec{F_{2}}\coord{4}{-1}$ (en kN, un carreau pour $1$~kN).
\begin{center}
\begin{tikzpicture}[x=0.45cm,y=0.45cm]
  \quadrillage{-0.5}{-2.5}{10.5}{3.5}
  \draw[vecteur] (0,0) -- node[midway,above left,font=\footnotesize]{$\vec{F_{1}}$} (3,2);
  \draw[vecteurB] (0,0) -- node[midway,below,font=\footnotesize]{$\vec{F_{2}}$} (4,-1);
  \pt{0,0}{left}{O}
\end{tikzpicture}
\end{center}
\begin{enumerate}[label=\alph*)]
  \item Calculer les coordonnées de la résultante $\vec{R}=\vec{F_{1}}+\vec{F_{2}}$.
  \item Tracer $\vec{R}$ sur la figure, en traçant $\vec{F_{2}}$ à la suite de
        $\vec{F_{1}}$.
  \item Un troisième remorqueur ajoute une force $\vec{F_{3}}$ pour que la résultante des
        trois forces soit $\coord{10}{0}$. Calculer les coordonnées de $\vec{F_{3}}$.
\end{enumerate}

\partiefiche{Parallélogramme et coordonnées d'un point}

\exo[Point défini par une égalité]
On donne $A(2\,;-1)$, $B(4\,;2)$, $C(0\,;5)$ et $\vec{u}\coord{-3}{4}$.
\begin{enumerate}[label=\alph*)]
  \item Calculer les coordonnées du point $E$ tel que $\vect{AE}=\vec{u}$.
  \item Calculer les coordonnées du point $D$ tel que $ABCD$ soit un parallélogramme.
  \item Calculer les coordonnées du point $F$ tel que $\vect{BF}=\vect{AC}$.
\end{enumerate}

\exo[Parallélogramme ou non]
Pour chaque quadrilatère, dire s'il est un parallélogramme, en justifiant.
\begin{enumerate}[label=\alph*)]
  \item $ABCD$, avec $A(-2\,;1)$, $B(1\,;3)$, $C(4\,;0)$ et $D(1\,;-2)$.
  \item $EFGH$, avec $E(0\,;0)$, $F(2\,;3)$, $G(5\,;4)$ et $H(3\,;2)$.
\end{enumerate}

\exo[Un portique]
Le cadre d'un portique de levage a la forme d'un parallélogramme $ABCD$. Dans un repère
en décimètres, trois fixations sont placées en $A(0\,;0)$, $B(6\,;1)$ et $C(8\,;5)$.
\begin{enumerate}[label=\alph*)]
  \item Placer les points $A$, $B$ et $C$ dans un repère (un carreau pour $1$~dm).
  \item Calculer les coordonnées de la quatrième fixation $D$, puis placer $D$ et
        vérifier sur la figure.
  \item Un renfort relie $C$ à un point $G$ tel que $ABGC$ soit aussi un
        parallélogramme. Calculer les coordonnées de $G$, puis placer $G$.
\end{enumerate}

\partiefiche{Colinéarité de deux vecteurs}

\exo[Colinéaires ou non]
Pour chaque couple de vecteurs, dire s'ils sont colinéaires : on cherche d'abord un
multiple évident, sinon on calcule le déterminant.
\begin{enumerate}[label=\alph*)]
  \item $\vec{u}\coord{3}{-1}$ et $\vec{v}\coord{-9}{3}$
  \item $\vec{u}\coord{2}{5}$ et $\vec{v}\coord{4}{11}$
  \item $\vec{u}\coord{6}{4}$ et $\vec{v}\coord{9}{6}$
  \item $\vec{u}\coord{-4}{10}$ et $\vec{v}\coord{2}{-5}$
  \item $\vec{u}\coord{5}{3}$ et $\vec{v}\coord{8}{5}$
\end{enumerate}

\exo[Trouver une coordonnée]
\begin{enumerate}[label=\alph*)]
  \item Calculer le nombre réel $m$ pour que $\vec{u}\coord{4}{-6}$ et $\vec{v}\coord{m}{9}$
        soient colinéaires.
  \item Même question pour $\vec{a}\coord{3}{m}$ et $\vec{b}\coord{6}{8}$.
\end{enumerate}

\exo[Avec des phrases]
On donne $\vec{u}\coord{-2}{4}$.
\begin{enumerate}[label=\alph*)]
  \item Donner les coordonnées de deux vecteurs colinéaires à $\vec{u}$, l'un de même
        sens, l'autre de sens contraire.
  \item Le vecteur $\vec{w}\coord{3}{-7}$ est-il colinéaire à $\vec{u}$ ? Répondre par une
        phrase complète, justifiée par un calcul.
\end{enumerate}

\exo[Deux forces sur une pièce]
Une pièce mécanique subit deux forces $\vec{F_{1}}\coord{6}{-4}$ et
$\vec{F_{2}}\coord{-3}{2}$ (en newtons).
\begin{enumerate}[label=\alph*)]
  \item Les forces $\vec{F_{1}}$ et $\vec{F_{2}}$ ont-elles la même direction ? Le même
        sens ?
  \item Calculer les coordonnées de la résultante $\vec{R}=\vec{F_{1}}+\vec{F_{2}}$.
  \item La résultante $\vec{R}$ a-t-elle la même direction que $\vec{F_{1}}$ ?
\end{enumerate}

\partiefiche{Droites parallèles, points alignés}

\exo[Parallèles, alignés]
On donne $A(-1\,;-2)$, $B(1\,;2)$, $C(3\,;-1)$, $D(4\,;1)$, $E(3\,;6)$ et $F(0\,;1)$.
\begin{enumerate}[label=\alph*)]
  \item Les droites $(AB)$ et $(CD)$ sont-elles parallèles ?
  \item Les points $A$, $B$ et $E$ sont-ils alignés ?
  \item Les points $A$, $B$ et $F$ sont-ils alignés ?
\end{enumerate}

\exo[Droites données par une équation]
\rappel{la droite d'équation cartésienne $ax+by+c=0$ a pour vecteur directeur
$\vec{u}\coord{-b}{a}$.}
Pour chaque couple de droites, donner un vecteur directeur de chacune, puis dire si les
droites sont parallèles.
\begin{enumerate}[label=\alph*)]
  \item $(d_{1})$ : $2x-y+3=0$ et $(d_{2})$ : $4x-2y-1=0$
  \item $(d_{3})$ : $x+3y-2=0$ et $(d_{4})$ : $2x+5y=0$
  \item $(d_{5})$ : $3x-6y+1=0$ et $(d_{6})$ : $-x+2y+4=0$
\end{enumerate}

\exo[La trajectoire d'un robot]
Un robot d'entrepôt va en ligne droite de $A(1\,;2)$ à $B(5\,;4)$ (repère en mètres), puis
continue tout droit.
\begin{enumerate}[label=\alph*)]
  \item Passera-t-il par le point $C(11\,;7)$ ?
  \item Passera-t-il par le point $D(9\,;7)$ ?
\end{enumerate}
Répondre à chaque question par une phrase complète.

\exo[Un trapèze]
On donne $A(0\,;0)$, $B(6\,;2)$, $C(5\,;5)$ et $D(2\,;4)$.
\begin{enumerate}[label=\alph*)]
  \item Montrer que les droites $(AB)$ et $(DC)$ sont parallèles.
  \item Les droites $(AD)$ et $(BC)$ sont-elles parallèles ?
  \item En déduire la nature du quadrilatère $ABCD$.
\end{enumerate}

\partiefiche{Coordonnées du milieu d'un segment}

\exo[Calculer des milieux]
Pour chaque segment, calculer les coordonnées de son milieu.
\begin{enumerate}[label=\alph*)]
  \item $[AB]$, avec $A(2\,;5)$ et $B(6\,;-1)$
  \item $[CD]$, avec $C(-3\,;4)$ et $D(1\,;-2)$
  \item $[EF]$, avec $E(-5\,;-2)$ et $F(1\,;4)$
\end{enumerate}

\exo[Retrouver une extrémité]
\begin{enumerate}[label=\alph*)]
  \item $K(1\,;3)$ est le milieu de $[AB]$, avec $A(-2\,;5)$. Calculer les coordonnées
        de $B$.
  \item $M(0\,;-1)$ est le milieu de $[CD]$, avec $C(3\,;2)$. Calculer les coordonnées
        de $D$.
\end{enumerate}

\exo[Le centre d'un parallélogramme]
On donne $A(-1\,;1)$, $B(3\,;0)$ et $C(5\,;3)$.
\begin{enumerate}[label=\alph*)]
  \item Calculer les coordonnées du point $D$ tel que $ABCD$ soit un parallélogramme.
  \item Calculer les coordonnées du milieu de $[AC]$, puis du milieu de $[BD]$.
        Que constate-t-on ?
\end{enumerate}

\partiefiche{Norme d'un vecteur}

\exo[Normes et distances]
Le repère est orthonormé.
\begin{enumerate}[label=\alph*)]
  \item Calculer la norme des vecteurs $\vec{u}\coord{5}{12}$, $\vec{v}\coord{-3}{0}$ et
        $\vec{w}\coord{2}{-2}$.
  \item On donne $A(1\,;-2)$ et $B(4\,;2)$. Calculer la distance $AB$.
  \item On donne $C(-1\,;3)$ et $D(2\,;1)$. Calculer $CD$, puis en donner une valeur
        approchée au centième.
\end{enumerate}

\exo[Nature d'un triangle]
Le repère est orthonormé. On donne $A(1\,;0)$, $B(5\,;2)$ et $C(3\,;6)$.
\begin{enumerate}[label=\alph*)]
  \item Calculer les longueurs $AB$, $AC$ et $BC$.
  \item Montrer que le triangle $ABC$ est isocèle, puis qu'il est rectangle.
\end{enumerate}

\exo[Les haubans d'une antenne]
\begin{minipage}[t]{0.52\linewidth}\raggedright
Une antenne verticale $[OS]$ est tenue par deux câbles $[SA]$ et $[SB]$. Dans un repère
orthonormé en mètres, $O(0\,;0)$, $S(0\,;12)$, $A(5\,;0)$ et $B(-9\,;0)$.
Calculer la longueur de chaque câble, puis la longueur totale de câble.
\end{minipage}\hfill
\begin{minipage}[t]{0.44\linewidth}
\vspace{0pt}\centering
\begin{tikzpicture}[scale=0.2]
  \draw[fig] (-10,0) -- (6,0);
  \draw[line width=1.6pt,black!70] (0,0) -- (0,12);
  \draw[coulCourbe,line width=0.9pt] (0,12) -- (5,0);
  \draw[coulCourbe,line width=0.9pt] (0,12) -- (-9,0);
  \pt{0,0}{below}{O} \pt{0,12}{above}{S} \pt{5,0}{below}{A} \pt{-9,0}{below}{B}
\end{tikzpicture}
\end{minipage}

\partiefiche{Définition du produit scalaire}

\exo[Avec un angle]
Pour chaque cas, calculer $\vec{u}\cdot\vec{v}$, sachant que $\alpha$ est l'angle entre
$\vec{u}$ et $\vec{v}$.
\begin{enumerate}[label=\alph*)]
  \item $\norme{\vec{u}}=5$, $\norme{\vec{v}}=4$ et $\alpha=60^{\circ}$
  \item $\norme{\vec{u}}=3$, $\norme{\vec{v}}=6$ et $\alpha=120^{\circ}$
  \item $\norme{\vec{u}}=2$, $\norme{\vec{v}}=7$ et $\alpha=90^{\circ}$ ; que peut-on dire
        des vecteurs $\vec{u}$ et $\vec{v}$ ? (phrase complète)
  \item $\norme{\vec{u}}=4$, $\norme{\vec{v}}=5$ et $\alpha=35^{\circ}$ (au dixième)
\end{enumerate}

\exo[Dans des triangles]
\begin{enumerate}[label=\alph*)]
  \item $ABC$ est un triangle équilatéral de côté $5$. Calculer $\vect{AB}\cdot\vect{AC}$.
  \item Dans le triangle $EFG$, $EF=4$, $EG=6$ et $\widehat{FEG}=120^{\circ}$. Calculer
        $\vect{EF}\cdot\vect{EG}$.
  \item $RST$ est un triangle rectangle en $R$. Calculer $\vect{RS}\cdot\vect{RT}$.
\end{enumerate}

\exo[Le travail d'une force]
\rappel{le travail d'une force $\vec{F}$ d'intensité $F$ (en N), pour un déplacement de
longueur $d$ (en m) qui fait un angle $\alpha$ avec la force, est
$W=F\times d\times\cos(\alpha)$ (en J).}
\begin{enumerate}[label=\alph*)]
  \item Un enfant tire un traîneau sur $25$~m avec une force de $40$~N qui fait un angle de
        $35^{\circ}$ avec le déplacement (schéma ci-dessous). Calculer le travail de la
        force, au dixième.
  \begin{center}
  \solidetire{35}{$\vec{F}$}{$40$~N}{$25$~m}
  \end{center}
  \item Un remonte-pente tire un skieur sur $200$~m avec une force de $300$~N qui fait un
        angle de $20^{\circ}$ avec le déplacement. Faire un schéma comme celui du a), puis
        calculer le travail de la force, au dixième.
\end{enumerate}

\partiefiche{Produit scalaire et projeté orthogonal}

\exo[Dans un rectangle]
\begin{minipage}[t]{0.55\linewidth}\raggedright
$ABCD$ est un rectangle avec $AB=5$ et $AD=3$. Calculer :
\begin{enumerate}[label=\alph*)]
  \item $\vect{AB}\cdot\vect{AC}$ ;
  \item $\vect{AD}\cdot\vect{AC}$ ;
  \item $\vect{AB}\cdot\vect{AD}$.
\end{enumerate}
\end{minipage}\hfill
\begin{minipage}[t]{0.42\linewidth}
\vspace{0pt}\centering
\begin{tikzpicture}[scale=0.42]
  \draw[fig] (0,0) -- (5,0) -- (5,3) -- (0,3) -- cycle;
  \draw[aide] (0,0) -- (5,3);
  \pt{0,0}{below left}{A} \pt{5,0}{below right}{B} \pt{5,3}{above right}{C}
  \pt{0,3}{above left}{D}
\end{tikzpicture}
\end{minipage}

\exo[Dans un carré]
\begin{minipage}[t]{0.55\linewidth}\raggedright
$ABCD$ est un carré de côté $4$, de centre $O$ ; $I$ est le milieu de $[AB]$. Calculer :
\begin{enumerate}[label=\alph*)]
  \item $\vect{AB}\cdot\vect{AC}$ ;
  \item $\vect{AB}\cdot\vect{AO}$ ;
  \item $\vect{AB}\cdot\vect{AD}$ ;
  \item $\vect{AD}\cdot\vect{AO}$.
\end{enumerate}
\end{minipage}\hfill
\begin{minipage}[t]{0.42\linewidth}
\vspace{0pt}\centering
\begin{tikzpicture}[scale=0.5]
  \draw[fig] (0,0) -- (4,0) -- (4,4) -- (0,4) -- cycle;
  \draw[aide] (0,0) -- (4,4) (4,0) -- (0,4);
  \draw[aide] (2,2) -- (2,0);
  \pt{0,0}{below left}{A} \pt{4,0}{below right}{B} \pt{4,4}{above right}{C}
  \pt{0,4}{above left}{D} \pt{2,2}{above}{O} \pt{2,0}{below}{I}
\end{tikzpicture}
\end{minipage}

\exo[Avec le projeté]
Dans chaque triangle, $H$ est le projeté orthogonal du troisième sommet sur la droite
qui porte le premier vecteur.
\begin{center}
\begin{tikzpicture}[scale=0.45]
  \draw[aide] (4,3) -- (4,0);
  \draw[black!60] (4,0.35) -- (3.65,0.35) -- (3.65,0);
  \draw[fig] (0,0) -- (6,0) -- (4,3) -- cycle;
  \pt{0,0}{below left}{E} \pt{6,0}{below right}{F} \pt{4,3}{above}{G} \pt{4,0}{below}{H}
  \node[font=\footnotesize] at (3,-1.8) {$EF=6$ et $EH=4$};
\end{tikzpicture}
\hspace{0.6cm}
\begin{tikzpicture}[scale=0.45]
  \draw[aide] (-2,0) -- (0,0);
  \draw[aide] (-2,3) -- (-2,0);
  \draw[black!60] (-2,0.35) -- (-1.65,0.35) -- (-1.65,0);
  \draw[fig] (0,0) -- (3,0) -- (-2,3) -- cycle;
  \pt{0,0}{below}{K} \pt{3,0}{below right}{L} \pt{-2,3}{above}{M} \pt{-2,0}{below}{H}
  \node[font=\footnotesize] at (0.5,-1.8) {$KL=3$ et $KH=2$};
\end{tikzpicture}
\end{center}
Calculer $\vect{EF}\cdot\vect{EG}$, puis $\vect{KL}\cdot\vect{KM}$.

\exo[Projection d'un vecteur sur un axe]
\begin{enumerate}[label=\alph*)]
  \item Sur un axe dirigé par $\vec{u}$, avec $\norme{\vec{u}}=3$, le vecteur $\vec{v}$
        de norme $10$ fait un angle de $30^{\circ}$ avec $\vec{u}$. Calculer la projection
        $p$ de $\vec{v}$ sur l'axe, au centième, puis $\vec{u}\cdot\vec{v}$, au dixième.
  \item Même question avec $\norme{\vec{u}}=4$, $\norme{\vec{v}}=6$ et un angle de
        $120^{\circ}$.
\end{enumerate}
\begin{center}
\vecteuraxe{30}\hspace{0.4cm}\vecteuraxe{120}
\end{center}

\partiefiche{Produit scalaire dans un repère orthonormé}

\exo[Avec les coordonnées]
Le repère est orthonormé.
\begin{enumerate}[label=\alph*)]
  \item Calculer $\vec{u}\cdot\vec{v}$ pour $\vec{u}\coord{2}{5}$ et $\vec{v}\coord{3}{-1}$.
  \item Calculer $\vec{a}\cdot\vec{b}$ pour $\vec{a}\coord{-4}{2}$ et $\vec{b}\coord{1}{2}$.
        Que peut-on dire de ces vecteurs ? (phrase complète)
  \item On donne $A(0\,;1)$, $B(3\,;2)$ et $C(2\,;5)$. Calculer $\vect{AB}\cdot\vect{AC}$.
\end{enumerate}

\exo[Perpendiculaires ou non]
Le repère est orthonormé. Pour chaque cas, dire si les droites sont perpendiculaires.
\begin{enumerate}[label=\alph*)]
  \item $(AB)$ et $(CD)$, avec $A(1\,;2)$, $B(3\,;6)$, $C(0\,;5)$ et $D(4\,;3)$.
  \item $(EF)$ et $(GH)$, avec $E(-2\,;0)$, $F(1\,;1)$, $G(0\,;3)$ et $H(2\,;-2)$.
\end{enumerate}

\exo[Parallèles, perpendiculaires ou ni l'un ni l'autre]
Le repère est orthonormé. Pour chaque couple de droites, donner un vecteur directeur de
chacune (vecteur $\vec{u}\coord{-b}{a}$ pour $ax+by+c=0$), puis dire si les droites sont
parallèles, perpendiculaires, ou ni l'un ni l'autre.
\begin{enumerate}[label=\alph*)]
  \item $(d_{1})$ : $2x+y-1=0$ et $(d_{2})$ : $x-2y+3=0$
  \item $(d_{3})$ : $3x-y=0$ et $(d_{4})$ : $x+2y-5=0$
  \item $(d_{5})$ : $4x-3y+2=0$ et $(d_{6})$ : $3x+4y-1=0$
  \item $(d_{7})$ : $x-y=0$ et $(d_{8})$ : $2x-2y+5=0$
\end{enumerate}

\exo[Un triangle rectangle]
Le repère est orthonormé. On donne $K(-1\,;1)$, $L(3\,;3)$ et $M(4\,;1)$.
Le triangle $KLM$ est-il rectangle en $L$ ? Répondre par une phrase complète.

\exo[L'équerrage d'un cadre]
Un menuisier vérifie un cadre de porte $ABCD$. Dans un repère orthonormé en
centimètres, $A(0\,;0)$, $B(120\,;9)$ et $D(-6\,;80)$.
\begin{enumerate}[label=\alph*)]
  \item Calculer les coordonnées des vecteurs $\vect{AB}$ et $\vect{AD}$.
  \item Le cadre a-t-il un angle droit en $A$ ?
  \item Calculer les coordonnées du point $C$ tel que $ABCD$ soit un parallélogramme.
  \item Le menuisier mesure les deux diagonales. Calculer $AC$ et $BD$ au dixième ;
        que constate-t-on ?
\end{enumerate}

\partiefiche{Propriétés du produit scalaire}

\exo[Développer]
On donne $\norme{\vec{u}}=2$, $\norme{\vec{v}}=5$ et $\vec{u}\cdot\vec{v}=3$. Calculer :
\begin{enumerate}[label=\alph*)]
  \item $\vec{u}\cdot(\vec{u}+\vec{v})$ ;
  \item $(\vec{u}+\vec{v})\cdot(\vec{u}-\vec{v})$ ;
  \item $(3\vec{u})\cdot(-2\vec{v})$ ;
  \item $\norme{\vec{u}+\vec{v}}^{2}$, puis $\norme{\vec{u}+\vec{v}}$ au centième ;
  \item $\norme{\vec{u}-\vec{v}}^{2}$.
\end{enumerate}

\exo[Changer l'ordre, le sens]
On sait que $\vect{AB}\cdot\vect{AC}=8$. Calculer $\vect{AC}\cdot\vect{AB}$,
$\vect{AB}\cdot\vect{CA}$, $(2\vect{AB})\cdot\vect{AC}$ et $\vect{BA}\cdot(3\vect{AC})$.

\exo[Dans un rectangle]
$ABCD$ est un rectangle avec $AB=6$ et $AD=4$ ; $I$ est le milieu de $[DC]$.
\begin{enumerate}[label=\alph*)]
  \item Faire une figure.
  \item Avec la relation de Chasles, montrer que $\vect{AC}=\vect{AB}+\vect{AD}$ et que
        $\vect{AI}=\vect{AD}+\frac{1}{2}\vect{AB}$.
  \item Que valent $\vect{AB}\cdot\vect{AB}$, $\vect{AD}\cdot\vect{AD}$ et
        $\vect{AB}\cdot\vect{AD}$ ?
  \item En développant, calculer $\vect{AC}\cdot\vect{AI}$.
\end{enumerate}

\partiefiche{Travail d'une force}

\exo[Une caisse tirée]
Une caisse avance de $8$~m sur un sol horizontal. Elle est tirée par une corde
($\vec{F}$ : $300$~N, à $25^{\circ}$ du sol) ; elle subit son poids ($\vec{P}$ :
$500$~N) et le frottement du sol ($\vec{f}$ : $120$~N, opposé au déplacement).
\begin{center}
\troisforces{25}{$25^{\circ}$}{$\vec{F}$}{$\vec{P}$}{$\vec{f}$}{$8$~m}
\end{center}
Calculer, au dixième, le travail du poids, de la traction et du frottement, puis le
travail de la résultante des trois forces.

\exo[Moteur, résistant ou nul]
Une force de $100$~N s'exerce sur un objet qui se déplace de $5$~m. Pour chaque angle
$\alpha$ entre la force et le déplacement, calculer le travail au dixième, puis dire par
une phrase complète s'il est moteur, résistant ou nul.
\begin{enumerate}[label=\alph*)]
  \item $\alpha=30^{\circ}$
  \item $\alpha=90^{\circ}$
  \item $\alpha=150^{\circ}$
  \item $\alpha=180^{\circ}$
\end{enumerate}

\exo[Sur une piste en pente]
Un chariot de $700$~N est hissé sur $150$~m le long d'une piste inclinée de $20^{\circ}$,
par un câble ($\vec{T}$ : $400$~N, à $15^{\circ}$ de la piste). La réaction $\vec{R}$ de
la piste lui est perpendiculaire.
\begin{center}
\begin{tikzpicture}[scale=0.6]
  \draw[fig] (-1,0) -- (8,0);
  \draw (2,0) arc (0:20:2);
  \node[font=\footnotesize,anchor=west] at (2.05,0.35) {$20^{\circ}$};
  \begin{scope}[rotate=20]
    \draw[fig] (0,0) -- (8.5,0);
    \draw[fig,fill=black!6] (3.3,0) rectangle (4.7,0.9);
    \draw[axe] (2,0.45) -- (9,0.45);
    \draw[deplacement] (4,0.45) -- node[pos=0.55,below right,font=\footnotesize]{déplacement\,: $150$~m} (7.8,0.45);
    \draw[vecteur] (4,0.45) -- node[pos=0.75,above left,font=\footnotesize]{$\vec{T}$} ++({3*cos(15)},{3*sin(15)});
    \draw[vecteurC] (4,0.45) -- node[pos=0.8,left,font=\footnotesize]{$\vec{R}$} ++(0,2.2);
    \pt{4,0.45}{below left}{G}
  \end{scope}
  \draw[vecteurB] ({4*cos(20)-0.45*sin(20)},{4*sin(20)+0.45*cos(20)}) -- node[pos=0.8,left,font=\footnotesize]{$\vec{P}$} ++(0,-2.6);
\end{tikzpicture}
\end{center}
\begin{enumerate}[label=\alph*)]
  \item L'angle entre le poids $\vec{P}$ et le déplacement mesure $110^{\circ}$.
        Calculer le travail du poids, au dixième. Est-il moteur ou résistant ?
  \item Calculer le travail de la traction $\vec{T}$, au dixième.
  \item Que vaut le travail de la réaction $\vec{R}$ ? Justifier.
  \item En déduire le travail de la résultante des trois forces.
\end{enumerate}

\exo[Un wagonnet]
Un wagonnet avance de $50$~m sur un rail rectiligne, tiré par un câble qui exerce une
force de $800$~N dont la direction fait un angle de $15^{\circ}$ avec le rail. Faire le schéma (solide, centre de gravité, axe
du déplacement, force), puis calculer, au dixième, la projection de la force sur l'axe
du rail et le travail de la force.

\partiefiche{Calcul d'un angle}

\exo[Angle entre deux vecteurs]
Le repère est orthonormé. Pour chaque couple de vecteurs, calculer l'angle qu'ils
forment, au dixième de degré.
\begin{enumerate}[label=\alph*)]
  \item $\vec{u}\coord{2}{1}$ et $\vec{v}\coord{1}{3}$
  \item $\vec{a}\coord{3}{-1}$ et $\vec{b}\coord{1}{2}$
  \item $\vec{c}\coord{-2}{1}$ et $\vec{d}\coord{3}{1}$
\end{enumerate}

\exo[Les angles d'un triangle]
Le repère est orthonormé. On donne $A(0\,;0)$, $B(4\,;0)$ et $C(2\,;3)$.
\begin{enumerate}[label=\alph*)]
  \item Calculer l'angle $\widehat{BAC}$, au dixième de degré.
  \item Calculer de même l'angle $\widehat{ABC}$.
  \item En déduire l'angle $\widehat{ACB}$, puis la nature du triangle $ABC$.
\end{enumerate}

\exo[Une porte renforcée]
\begin{minipage}[t]{0.5\linewidth}\raggedright
Une porte rectangulaire $OEKA$ mesure $3$~m sur $2$~m. Pour la renforcer, on fixe deux
barres $[OB]$ et $[OD]$, avec $B(1\,;2)$ et $D(3\,;1{,}5)$ (repère orthonormé en
mètres, d'origine $O$). Le renfort est efficace si l'angle $\alpha=\widehat{BOD}$ est
compris entre $35^{\circ}$ et $45^{\circ}$.
\end{minipage}\hfill
\begin{minipage}[t]{0.46\linewidth}
\vspace{0pt}\centering
\begin{tikzpicture}[scale=1]
  \draw[fig,fill=black!4] (0,0) rectangle (3,2);
  \draw[coulCourbe,line width=1.1pt] (0,0) -- (1,2);
  \draw[coulCourbe,line width=1.1pt] (0,0) -- (3,1.5);
  \draw (0.35,0.175) arc (26.57:63.43:0.39);
  \node[font=\footnotesize] at (0.5,0.55) {$\alpha$};
  \pt{0,0}{below left}{O} \pt{3,0}{below right}{E} \pt{3,2}{above right}{K}
  \pt{0,2}{above left}{A} \pt{1,2}{above}{B} \pt{3,1.5}{right}{D}
\end{tikzpicture}
\end{minipage}
\begin{enumerate}[label=\alph*)]
  \item Calculer les coordonnées de $\vect{OB}$ et de $\vect{OD}$, puis les longueurs $OB$
        et $OD$, au centième.
  \item Calculer $\vect{OB}\cdot\vect{OD}$.
  \item Calculer l'angle $\alpha$, au dixième de degré. Le renfort est-il efficace ?
\end{enumerate}

\partiefiche{Théorème d'Al-Kashi}

\exo[Calculer un côté]
\rappel{théorème d'Al-Kashi : dans un triangle $ABC$,
$BC^{2}=AB^{2}+AC^{2}-2\times AB\times AC\times\cos\big(\widehat{BAC}\big)$.}
Pour chaque triangle $ABC$, calculer $BC$ (au centième si besoin).
\begin{enumerate}[label=\alph*)]
  \item $AB=3$, $AC=5$ et $\widehat{BAC}=60^{\circ}$
  \item $AB=5$, $AC=3$ et $\widehat{BAC}=120^{\circ}$
  \item $AB=7$, $AC=10$ et $\widehat{BAC}=40^{\circ}$
\end{enumerate}

\exo[Un bras de robot]
\begin{minipage}[t]{0.52\linewidth}\raggedright
Le bras d'un robot est formé de deux segments articulés en $C$ : $[SC]$ de $50$~cm et
$[CP]$ de $40$~cm. Calculer la distance $SP$ entre l'épaule $S$ et la pince $P$, au
dixième, quand l'angle $\widehat{SCP}$ mesure $100^{\circ}$, puis quand il mesure
$60^{\circ}$.
\end{minipage}\hfill
\begin{minipage}[t]{0.44\linewidth}
\vspace{0pt}\centering
\begin{tikzpicture}[scale=0.55]
  \draw[line width=1.8pt,black!65] (0,0) -- (5,0) -- (5.69,3.94);
  \draw[aide] (0,0) -- (5.69,3.94);
  \draw (4.3,0) arc (180:80:0.7);
  \node[font=\footnotesize] at (4.1,0.75) {$100^{\circ}$};
  \pt{0,0}{below left}{S} \pt{5,0}{below right}{C} \pt{5.69,3.94}{above}{P}
\end{tikzpicture}
\end{minipage}

\exo[Deux bateaux]
\begin{minipage}[t]{0.52\linewidth}\raggedright
Deux bateaux partent en même temps du port $O$, en ligne droite, sur deux routes qui font
entre elles un angle de $50^{\circ}$. Une heure plus tard, le premier est en $A$, à
$12$~km du port, et le second en $B$, à $9$~km du port. Calculer la distance $AB$ entre
les deux bateaux, au centième.
\end{minipage}\hfill
\begin{minipage}[t]{0.44\linewidth}
\vspace{0pt}\centering
\begin{tikzpicture}[scale=0.28]
  \draw[fig] (0,0) -- node[below,font=\footnotesize]{$12$~km} (12,0);
  \draw[fig] (0,0) -- node[above left,font=\footnotesize]{$9$~km} ({9*cos(50)},{9*sin(50)});
  \draw[aide] (12,0) -- ({9*cos(50)},{9*sin(50)});
  \draw (2.5,0) arc (0:50:2.5);
  \node[font=\footnotesize] at (4.1,1.7) {$50^{\circ}$};
  \pt{0,0}{below left}{O} \pt{12,0}{below right}{A} \pt{{9*cos(50)},{9*sin(50)}}{above}{B}
\end{tikzpicture}
\end{minipage}

\end{multicols*}

\end{document}
